%=============================================================================
% Wave 4, Iteration 13: Materials Science Update
% AGENT: MATERIALS SCIENCE (W4i13) | Model: Opus 4.6
%
% INHERITED FROM W4i10 MatSci:
%   - Si3N4 MEMS: Q>=10^9, 1.2 ng, 620 kHz, threshold 2.2e8 (4.5x margin)
%   - Nb Q0 5e10 baseline, 2e11 target; SQMS $125M renewal
%   - He-3 KILLED (34 orders gap)
%   - Dark SRF INVALIDATED (W4i12) --- passive detection only
%   - Passive >> active by 14 orders (Passive Superiority Theorem)
%   - Three-track roadmap (Nb/Si3N4/He-3) established
%   - Supply chain: Norcada, Atomica, university cleanrooms
%   - Cryo: Bluefors XLD400sl, $0.8-1.5M
%
% W4i13 MANDATE:
%   1. Go BEYOND standard MatSci: novel materials for passive FOM
%   2. Crystalline resonators (Si, SiC) vs Si3N4: Q>10^10 implications
%   3. Levitated nanoparticle optomechanics for psi-force sensing
%   4. AlN piezoelectric PnC membranes: new competitor?
%   5. Nb3Sn and alternative SRF materials for Phase I/III
%   6. Diamond NV centers as psi-field sensors
%   7. Photonic crystal cavities for optomechanical psi readout
%   8. O-doping as alternative to N-doping for Nb SRF
%
% Date: 20 March 2026
%=============================================================================

\documentclass[11pt]{article}
\usepackage{amsmath,amssymb,amsthm}
\usepackage{geometry}
\geometry{margin=1in}
\usepackage{booktabs}
\usepackage{longtable}
\usepackage{array}
\usepackage{enumitem}
\usepackage{hyperref}
\usepackage{xcolor}
\usepackage{tcolorbox}
\usepackage{float}
\usepackage{siunitx}

\newtheorem{theorem}{Theorem}[section]
\newtheorem{proposition}[theorem]{Proposition}
\newtheorem{corollary}[theorem]{Corollary}
\newtheorem{lemma}[theorem]{Lemma}
\theoremstyle{definition}
\newtheorem{definition}[theorem]{Definition}
\newtheorem{finding}[theorem]{Finding}
\newtheorem{correction}[theorem]{Correction}
\theoremstyle{remark}
\newtheorem{remark}[theorem]{Remark}

\newcommand{\ee}[1]{\times 10^{#1}}
\newcommand{\Meff}{M_{\rm eff}}
\newcommand{\Qpsi}{Q_\psi}
\newcommand{\dpsi}{\delta\psi_{\rm EM}}
\newcommand{\psigrav}{\psi_{\rm grav}}
\newcommand{\psiEM}{\delta\psi_{\rm EM}}
\newcommand{\lambdaone}{|\lambda-1|}
\newcommand{\lambdabound}{1.56\ee{-9}}
\newcommand{\Tc}{T_{\rm c}}
\newcommand{\lambdaL}{\lambda_{\rm L}}
\newcommand{\FOMpassive}{\mathrm{FOM}_{\rm passive}}
\newcommand{\FOMMEMS}{\mathrm{FOM}_{\rm MEMS}}
\newcommand{\lam}{|\lambda-1|}
\newcommand{\Qmech}{Q_{\rm mech}}
\newcommand{\SNR}{\mathrm{SNR}}
\newcommand{\confirmed}[1]{\textcolor{blue}{\textbf{CONFIRMED:} #1}}
\newcommand{\corrected}[1]{\textcolor{red}{\textbf{CORRECTED:} #1}}
\newcommand{\caution}[1]{\textcolor{orange}{\textbf{CAUTION:} #1}}
\newcommand{\honest}[1]{\textcolor{violet}{\textbf{HONEST ASSESSMENT:} #1}}
\newcommand{\killed}[1]{\textcolor{red!70!black}{\textbf{KILLED:} #1}}
\newcommand{\verdict}[1]{\textcolor{green!50!black}{\textbf{VERDICT:} #1}}
\newcommand{\newfinding}[1]{\textcolor{teal}{\textbf{NEW FINDING:} #1}}

\title{Wave 4 Iteration 13: Materials Science Update\\
\large Crystalline Resonators Beyond Si$_3$N$_4$,\\
Levitated Nanoparticle Optomechanics,\\
AlN Piezoelectric PnC Membranes,\\
Nb$_3$Sn Alternative SRF,\\
and Novel Passive-Detection Material Survey}
\author{DFD Research Programme --- Materials Science Agent, W4i13}
\date{20 March 2026}

\begin{document}
\maketitle
\tableofcontents
\newpage

%=============================================================================
\section{Executive Summary: Top 5 Findings}
\label{sec:top5}
%=============================================================================

\begin{tcolorbox}[colback=blue!5!white, colframe=blue!50!black,
  title=\textbf{W4i13 TOP 5 FINDINGS --- Ranked by Programme Impact}]
\begin{enumerate}

\item \textbf{FINDING 1: Crystalline silicon soft-clamped resonators achieve
$Q > 10^{10}$ at 7~K --- a $10\times$ improvement over Si$_3$N$_4$.}
Ghadimi et al.\ (Nature Physics, 2022) demonstrated strained crystalline
silicon nanostrings in phononic crystals with $Q > 10^{10}$ at MHz
frequencies and 7~K. This improves the DFD MEMS $\Qpsi$ margin from
$4.5\times$ (at Si$_3$N$_4$ $Q = 10^9$) to $\mathbf{45\times}$ ---
a decisive margin for the in-gap diagnostic. Derivation chain:
$\Qpsi$ threshold $= 2.2\ee{8}$ (W4i11);
$Q_{\rm Si,cryo}/\Qpsi = 10^{10}/2.2\ee{8} = 45$.
\textbf{Impact: Phase~II MEMS detection becomes robust against
up to $45\times$ systematic degradation.} However, cryogenic operation
(7~K) is required, adding dilution refrigerator cost. Net cost impact:
marginal, since Phase~I already requires a dilution refrigerator.
Trade-off: crystalline Si requires cryogenic operation but delivers
order-of-magnitude $Q$ improvement; Si$_3$N$_4$ operates at room
temperature with $4.5\times$ margin. \textbf{Recommendation: dual-track ---
Si$_3$N$_4$ for room-temperature prototype, crystalline Si for cryogenic
high-confidence detection.}

\item \textbf{FINDING 2: Levitated nanoparticle optomechanics achieves
$\sim$166~yN force resolution --- a fundamentally different detection
geometry for psi-force sensing.}
Optically levitated silica nanospheres ($m \sim 10^{-18}$~kg,
$Q > 10^9$) in ultra-high vacuum have demonstrated force sensitivity
of $4.34$~zN$/\sqrt{\rm Hz}$ and force resolution of
$\sim 166$~yN ($1.7\ee{-22}$~N). This is a \textit{qualitatively
different} detection architecture from clamped MEMS: the nanoparticle
is mechanically isolated from all substrates, eliminating clamping
loss entirely. \textbf{For DFD:} The psi-force on a mass $m$ in
a gradient $\nabla\psi_{\rm grav}$ from a source mass $M$ at distance
$r$ is $F_\psi = (mc^2/2)|\nabla\psi| = mGM/r^2$ (Newtonian limit,
$\mu \to 1$). For $m = 10^{-18}$~kg, $M = 100$~kg, $r = 0.1$~m:
$F_\psi = 10^{-18} \times 6.67\ee{-11} \times 100/(0.01)
= 6.67\ee{-25}$~N.
This is $4000\times$ below yoctonewton resolution --- \textbf{still
infeasible for gravitational psi-force detection in the lab.}
However, for an EM-sourced $\dpsi$ perturbation at the SQMS bound
$\lam < 1.56\ee{-9}$, the force is further suppressed by $\lam^2$,
making it hopeless.
\verdict{Levitated optomechanics is KILLED for direct psi-force
sensing (4000$\times$ gap for gravitational, much worse for EM).
Retain as technology watch for future sub-yN sensitivity.}

\item \textbf{FINDING 3: Nb$_3$Sn achieves $Q_0 \geq 5\ee{10}$ at 2~K
and $\geq 2\ee{10}$ at 4~K --- enabling Phase~III at 4~K instead
of 20~mK.}
Nb$_3$Sn coated cavities (Fermilab/Cornell) demonstrate $Q_0$ values
competitive with N-doped Nb but at \textit{dramatically higher}
operating temperatures: $T_c = 18$~K vs 9.3~K for Nb. The cryogenic
cost reduction is transformative for Phase~III (256-cavity array):
replacing dilution refrigerators (\$1--2M each) with 4~K cryocoolers
(\$50--100K each) reduces the cryo budget by $\sim 20\times$. At 4~K,
the BCS surface resistance of Nb$_3$Sn scales as
$R_{\rm BCS} \propto (T/T_c)^2 \exp(-\Delta/k_BT)$
with $\Delta_{\rm Nb_3Sn}/k_B \approx 38$~K, giving
$R_{\rm BCS}(4~{\rm K}) \approx R_{\rm Nb}(1.5~{\rm K})$.
\textbf{Derivation of Phase~III impact:}
Phase~III requires 256 cavities at $Q_0 > 10^{10}$.
Nb at 20~mK: 256 dilution refrigerators at \$1M each $=$ \$256M cryo.
Nb$_3$Sn at 4~K: 256 cryocoolers at \$75K each $=$ \$19M cryo.
Saving: \$237M, reducing Phase~III total from \$175M+ to \$\sim 100M.
\caution{Nb$_3$Sn surface quality remains a challenge: residual
resistance from grain boundaries and surface oxide limits $Q_0$ at
low fields. Phase~I should use N-doped Nb (proven); Nb$_3$Sn is a
Phase~III upgrade path.}
Additionally, March 2026 research (SQMS-supported) shows oxygen doping
of Nb achieves surface resistance reductions comparable to nitrogen
doping, providing a second optimisation axis for Phase~I cavity
preparation.

\item \textbf{FINDING 4: AlN piezoelectric phononic crystal membranes
achieve $Q_m \times f = 1.5\ee{13}$~Hz at room temperature with
built-in piezoelectric actuation.}
Wieczorek et al.\ (APL, June 2025) demonstrated a 90~nm strained
AlN membrane phononic crystal with a 1.8~MHz defect mode achieving
$Q_m \times f = 1.5\ee{13}$~Hz at room temperature. For comparison,
Si$_3$N$_4$ achieves $Q \times f \sim 6\ee{14}$~Hz (Tsaturyan et al.\
2017). The AlN $Q \times f$ product is $40\times$ lower than best
Si$_3$N$_4$, \textbf{so AlN does not displace Si$_3$N$_4$ for
passive $\Qpsi$ detection.}
\textbf{However}, AlN has a unique advantage: built-in piezoelectricity
enables direct electrical actuation and readout without optical
interferometry, and direct coupling to superconducting microwave
circuits (transmon qubits). For the DFD programme, this means:
\begin{itemize}[nosep]
\item \textit{Phase~II hybrid readout}: AlN membrane coupled to
  a transmon qubit for quantum-limited displacement sensing,
  eliminating the optical interferometry chain.
\item \textit{Tunable phononic circuits}: In-situ frequency tuning
  of defect modes via DC bias on AlN electrodes.
\end{itemize}
\verdict{AlN does not replace Si$_3$N$_4$ for passive FOM but
opens a new readout pathway. Track as Phase~II readout option;
do not change the primary MEMS material selection.}

\item \textbf{FINDING 5: Monocrystalline 4H-SiC nanomechanical resonators
reach $Q = 2\ee{7}$ at room temperature with path to $10^{10}$ via
dissipation dilution.}
Sementilli et al.\ (Nano Letters, 2025) demonstrated monocrystalline
4H-SiC resonators with damping as low as 2.7~mHz and $Q = 2\ee{7}$
at room temperature --- the material limit for SiC --- an
order-of-magnitude improvement over prior crystalline SiC devices.
Crucially, the low \textit{linear} dissipation was achieved at
\textit{low stress}, making these devices mechanically robust.
The authors project that dissipation dilution and strain engineering
(soft clamping in phononic crystals) should yield $Q \sim 10^{10}$
at room temperature in SiC.

\textbf{DFD relevance:} SiC combines:
\begin{itemize}[nosep]
\item Optically active defect centres (silicon vacancies, V$_{\rm Si}$)
  for integrated spin-mechanical readout
\item Wide bandgap (3.26~eV) for low optical absorption
\item Piezoelectric coupling (weaker than AlN but nonzero in 4H polytype)
\item Chemical inertness and radiation hardness
\end{itemize}

If SiC reaches $Q \sim 10^{10}$ at room temperature, it would match
cryogenic crystalline Si without requiring a dilution refrigerator.
\textbf{This is speculative but high-impact: a room-temperature
$10^{10}$ resonator would give a $\Qpsi$ margin of $45\times$ without
cryogenics.}

\caution{$Q = 10^{10}$ in SiC is \textit{projected}, not demonstrated.
Current best is $2\ee{7}$. The gap to $10^{10}$ requires $500\times$
improvement from dissipation dilution, which is plausible based on
the Si$_3$N$_4$ precedent ($10^4 \to 10^9$ from bulk to soft-clamped)
but unproven in SiC.}

\verdict{SiC is a long-term watch item. Add to materials tracking
alongside Si$_3$N$_4$ and crystalline Si. No change to near-term
roadmap.}
\end{enumerate}
\end{tcolorbox}


%=============================================================================
\section{Finding 1: Crystalline Silicon Soft-Clamped Resonators}
\label{sec:crystalline_si}
%=============================================================================

\subsection{State of the Art}

The landmark result is strained crystalline silicon nanostrings
embedded in phononic crystal shields, achieving $Q > 10^{10}$ at
7~K with MHz mechanical frequencies (Ghadimi et al., Nature Physics
2022). The mechanism combines three effects:

\begin{enumerate}[nosep]
\item \textbf{Dissipation dilution}: Tensile pre-stress ($\sigma
  \sim 1$~GPa) stores elastic energy predominantly in the lossless
  stress field rather than in lossy bending deformation. The dilution
  factor is $D_Q = \sigma L^2/(E t^2)$, where $L$ is the string
  length, $E$ is Young's modulus, and $t$ is thickness.

\item \textbf{Soft clamping}: The phononic crystal provides an
  exponentially decaying mode profile at the boundaries, eliminating
  the sharp curvature kinks of hard-clamped resonators. This removes
  the dominant loss mechanism (clamping radiation) in conventional
  MEMS.

\item \textbf{Crystalline perfection}: Unlike amorphous Si$_3$N$_4$,
  crystalline Si has extremely low intrinsic material loss:
  $Q_{\rm int}^{-1} \sim 10^{-11}$ at cryogenic temperatures
  (limited by phonon-phonon Akhiezer damping). The material loss
  floor of Si$_3$N$_4$ is $Q_{\rm int}^{-1} \sim 10^{-8}$ (TLS
  losses in amorphous material).
\end{enumerate}

\subsection{DFD Passive FOM Comparison}

The passive MEMS figure of merit for DFD in-gap detection is:
\begin{equation}
\FOMMEMS = \frac{\Qmech}{m_{\rm eff} \cdot \omega_m^2},
\label{eq:fom_mems}
\end{equation}
where $\Qmech$ is the mechanical quality factor, $m_{\rm eff}$ is
the effective motional mass, and $\omega_m = 2\pi f_0$ is the angular
mechanical frequency. The psi-field induced displacement is:
\begin{equation}
x_\psi = \frac{F_\psi}{m_{\rm eff} \omega_m^2 / \Qmech}
= \frac{F_\psi \cdot \Qmech}{m_{\rm eff} \omega_m^2},
\label{eq:x_psi}
\end{equation}
so larger FOM means larger displacement for a given psi-force.

\begin{table}[H]
\centering
\caption{Passive MEMS FOM comparison: three material platforms.}
\label{tab:fom_compare}
\renewcommand{\arraystretch}{1.3}
\begin{tabular}{lccccc}
\toprule
\textbf{Platform} & $Q$ & $m_{\rm eff}$ (ng) & $f_0$ (MHz) &
  $T_{\rm op}$ & \textbf{$Q/Q_\psi^{\rm threshold}$} \\
\midrule
Si$_3$N$_4$ PnC (W3i7 Design~C) &
  $10^9$ & 1.2 & 0.62 & 300~K & $4.5\times$ \\
Crystalline Si PnC nanostring &
  $10^{10}$ & $\sim 1$--10 & $\sim 1$ & 7~K & $45\times$ \\
4H-SiC (projected, soft-clamped) &
  $\sim 10^{10}$ (projected) & TBD & $\sim 1$ & 300~K & $45\times$ (if achieved) \\
\bottomrule
\end{tabular}
\end{table}

\subsection{Centimetre-Scale Si$_3$N$_4$ Resonators}

A parallel development: centimetre-scale Si$_3$N$_4$ trampoline
resonators have achieved $Q = 6.6\ee{9}$ at 214~kHz at room
temperature (Bereyhi et al., Nature Communications 2024). This is
the highest $Q$ measured for any clamped resonator at room temperature.
The $\Qpsi$ margin at this $Q$ is $6.6\ee{9}/2.2\ee{8} = 30\times$,
approaching crystalline Si performance without cryogenics.

\begin{finding}[Crystalline Si vs Si$_3$N$_4$: Dual-Track Recommendation]
\label{find:dual_track}
\textbf{Near-term (Phase~II, 2027--28)}: Si$_3$N$_4$ at room temperature.
$Q = 10^9$ with $4.5\times$ margin. Fabrication chain mature (Norcada,
Delft, Yale). Cost: \$200--500K.

\textbf{High-confidence (Phase~II+, 2028--29)}: Crystalline Si at
7~K. $Q = 10^{10}$ with $45\times$ margin. Requires collaboration
with EPFL (Kippenberg group) or DTU Nanolab. Cost increment:
\$100--200K for device development (dilution refrigerator already
procured for Phase~I SRF).

\confirmed{Crystalline Si PnC resonators improve the in-gap
diagnostic margin by $10\times$ over Si$_3$N$_4$, at the cost of
requiring 7~K cryogenic operation. Since Phase~I already requires
a dilution refrigerator, the marginal cost of adding a Si MEMS
measurement is small (\$100--200K for device procurement and
mounting).}
\end{finding}


%=============================================================================
\section{Finding 2: Levitated Nanoparticle Optomechanics}
\label{sec:levitated}
%=============================================================================

\subsection{Current Performance}

The state of the art in levitated optomechanics (2024--2025):

\begin{itemize}[nosep]
\item Force sensitivity: $4.34$~zN$/\sqrt{\rm Hz}$ (parametric
  feedback cooling at $7.9\ee{-9}$~mbar)
\item Force resolution: $166 \pm 55$~yN (integration time 2751~s)
\item Mechanical $Q > 10^9$ (limited by residual gas pressure)
\item Ground state cooling achieved in 1D and 2D
\item Mass: typically $m \sim 10^{-19}$--$10^{-17}$~kg (silica
  nanospheres 50--500~nm diameter)
\end{itemize}

\subsection{DFD psi-Force Analysis}

The psi-force on a test mass $m$ from a gravitational source mass
$M$ at distance $r$ in the Newtonian limit ($\mu \to 1$):
\begin{equation}
F_\psi = \frac{mc^2}{2}|\nabla\psi_{\rm grav}|
= \frac{mGM}{r^2}.
\label{eq:f_psi_grav}
\end{equation}

For a levitated nanoparticle ($m = 10^{-18}$~kg) near a laboratory
source ($M = 100$~kg, $r = 0.1$~m):
\begin{equation}
F_\psi = \frac{10^{-18} \times 6.674\ee{-11} \times 10^2}{10^{-2}}
= 6.67\ee{-25}~\text{N} = 0.667~\text{yN}.
\label{eq:f_psi_num}
\end{equation}

This is $\sim 250\times$ below the best demonstrated force resolution
of $\sim 166$~yN. With integration time scaling as $t \propto
(\SNR)^2$, reaching $\SNR = 1$ at 0.667~yN would require:
\begin{equation}
t_{\rm req} = 2751~{\rm s} \times \left(\frac{166}{0.667}\right)^2
\approx 1.7\ee{8}~{\rm s} \approx 5.4~{\rm years}.
\label{eq:t_req}
\end{equation}

This is \textit{technically feasible} for a single continuous
measurement, but impractical due to drift, laser instability, and
vacuum degradation over such timescales.

\subsection{EM-Sourced psi-Force (Even Worse)}

For the EM-sourced perturbation $\dpsi$ detectable at the SQMS
bound $\lam < 1.56\ee{-9}$, the psi-force is further suppressed:
\begin{equation}
F_{\dpsi} \sim \lam \cdot \frac{m \cdot u_{\rm EM}}{M_{\rm eff}}
\sim 10^{-9} \times \frac{10^{-18} \times 10^4}{3\ee{6}}
\sim 10^{-29}~\text{N},
\label{eq:f_dpsi}
\end{equation}
which is $\sim 10^7 \times$ below yoctonewton resolution.

\begin{finding}[Levitated Nanoparticle psi-Force Sensing]
\label{find:levitated}
\killed{Direct psi-force sensing with levitated nanoparticles:
gravitational $F_\psi \sim 0.7$~yN vs resolution $\sim 166$~yN
($250\times$ gap, or 5.4~yr integration). EM-sourced $\dpsi$ force
$\sim 10^{-29}$~N ($10^7\times$ gap). Both channels are infeasible.}

\honest{Levitated optomechanics is the \textit{closest} any
non-DFD-specific technology comes to gravitational psi-force
detection (only $250\times$ gap vs $34$ orders for He-3). With
projected $10\times$ sensitivity improvements over 5--10~years,
the gravitational channel could become marginally feasible
($25\times$ gap $\to$ requiring $\sim 50$~day integration).
This is a long-term watch item.}
\end{finding}


%=============================================================================
\section{Finding 3: Nb$_3$Sn for Phase~III SRF Cavities}
\label{sec:nb3sn}
%=============================================================================

\subsection{Performance at Elevated Temperatures}

Nb$_3$Sn is an A15 compound superconductor with:
\begin{itemize}[nosep]
\item $T_c = 18.3$~K (vs 9.25~K for Nb)
\item $\Delta_{\rm Nb_3Sn}/k_B \approx 38$~K (vs 17.1~K for Nb)
\item $H_{c2} > 25$~T (vs $\sim 0.2$~T for Nb)
\item Lower critical field $H_{c1} \approx 50$~mT
\end{itemize}

Recent results from Fermilab/Cornell/JLab:
\begin{table}[H]
\centering
\caption{Nb$_3$Sn SRF cavity performance.}
\renewcommand{\arraystretch}{1.3}
\begin{tabular}{lccc}
\toprule
\textbf{Temperature} & $Q_0$ & $E_{\rm acc}$ & \textbf{Notes} \\
\midrule
2~K & $\geq 5\ee{10}$ & $\leq 15$~MV/m & Fermilab, vapor diffusion \\
4~K & $1$--$2\ee{10}$ & $\leq 15$~MV/m & 10--100$\times$ Nb at 4~K \\
4.4~K & $\sim 10^{10}$ & low--medium & Cryocooler-compatible \\
\bottomrule
\end{tabular}
\end{table}

\subsection{BCS Surface Resistance Analysis}

The BCS surface resistance scales as:
\begin{equation}
R_{\rm BCS}(T, f) = A \cdot \frac{f^2}{T}
  \exp\!\left(-\frac{\Delta}{k_B T}\right),
\label{eq:bcs_rs}
\end{equation}
where $A$ is a material-dependent prefactor. The ratio of Nb$_3$Sn
at 4~K to Nb at 2~K:
\begin{equation}
\frac{R_{\rm BCS}^{\rm Nb_3Sn}(4~{\rm K})}
     {R_{\rm BCS}^{\rm Nb}(2~{\rm K})}
= \frac{A_{\rm Nb_3Sn}}{A_{\rm Nb}} \cdot \frac{2}{4}
  \cdot \exp\!\left(-\frac{38}{4} + \frac{17.1}{2}\right)
= \frac{A_{\rm Nb_3Sn}}{A_{\rm Nb}} \cdot 0.5
  \cdot \exp(-0.95).
\label{eq:bcs_ratio}
\end{equation}

With $A_{\rm Nb_3Sn}/A_{\rm Nb} \sim 2$--3 (larger coherence
length $\xi$ in Nb$_3$Sn), this gives $R_{\rm BCS}$ ratio $\sim$
0.4--0.6. \textbf{Nb$_3$Sn at 4~K is roughly BCS-equivalent to
Nb at 1.5--2~K}, confirming the experimental $Q_0$ data.

\subsection{Phase~III Cost Impact}

\begin{finding}[Nb$_3$Sn Phase~III Upgrade Path]
\label{find:nb3sn_phase3}
\textbf{Phase~III specification}: 256-cavity phased array, each
cavity $Q_0 > 10^{10}$, total reach $\lam \sim 10^{-12}$.

With N-doped Nb at 20~mK:
\begin{itemize}[nosep]
\item 256 dilution refrigerators or large-scale DR system:
  \$100--256M
\item Power: 256 $\times$ 50~kW = 12.8~MW wall-plug (prohibitive
  if individual DRs; centralized system \$50--100M)
\end{itemize}

With Nb$_3$Sn at 4~K:
\begin{itemize}[nosep]
\item 256 Gifford-McMahon or pulse-tube cryocoolers at \$50--100K
  each: \$13--26M
\item Power: 256 $\times$ 10~kW = 2.6~MW wall-plug (feasible at
  a single facility)
\end{itemize}

\textbf{Saving: \$75--230M in cryogenic costs alone.}

\caution{Nb$_3$Sn surface preparation is not yet at the
reproducibility level of N-doped Nb. Residual resistance from grain
boundaries, tin segregation, and surface oxide remains a limiting
factor. Electropolishing and oxypolishing of Nb$_3$Sn are active
research areas (JLab, 2025--2026). Phase~I must use Nb; Nb$_3$Sn is
a Phase~III-ready technology.}

\newfinding{Oxygen doping of Nb (SQMS, March 2026) achieves
comparable $R_s$ reduction to nitrogen doping. This provides a second
optimisation pathway for Phase~I: N-doped cavities as primary,
O-doped as backup/comparison.}
\end{finding}


%=============================================================================
\section{Finding 4: AlN Piezoelectric Phononic Crystal Membranes}
\label{sec:aln}
%=============================================================================

\subsection{Recent Results}

Wieczorek et al.\ (APL, June 2025; arXiv:2501.19151) demonstrated:
\begin{itemize}[nosep]
\item 90~nm strained AlN membrane with honeycomb phononic crystal
\item Localized defect mode at $f_0 = 1.8$~MHz
\item $Q_m \times f = 1.5\ee{13}$~Hz at room temperature
\item Built-in piezoelectric coefficient $d_{33} \sim 5$~pm/V
\end{itemize}

\subsection{Comparison with Si$_3$N$_4$}

\begin{table}[H]
\centering
\caption{AlN vs Si$_3$N$_4$ PnC membrane comparison.}
\renewcommand{\arraystretch}{1.3}
\begin{tabular}{lcc}
\toprule
\textbf{Property} & \textbf{AlN PnC} & \textbf{Si$_3$N$_4$ PnC} \\
\midrule
$Q \times f$ product & $1.5\ee{13}$~Hz & $6\ee{14}$~Hz \\
$Q$ at $\sim 1$~MHz & $\sim 8\ee{6}$ & $\sim 10^9$ \\
Piezoelectric & Yes ($d_{33} \sim 5$~pm/V) & No \\
Optical readout & Possible & Standard \\
Microwave coupling & Direct (piezo) & Via capacitive transduction \\
Operating $T$ & 300~K & 300~K \\
Fabrication & Sputtered AlN on Si & LPCVD stoichiometric \\
\bottomrule
\end{tabular}
\end{table}

\subsection{DFD Programme Relevance}

The $Q$ deficit of AlN ($\sim 120\times$ below Si$_3$N$_4$) makes it
unsuitable as the primary passive detection element. However, the
piezoelectric coupling enables a qualitatively new readout scheme:

\begin{equation}
V_{\rm piezo} = \frac{d_{33} \cdot x_\psi}{t_{\rm AlN}}
\cdot \frac{1}{C_{\rm parasitic}},
\label{eq:v_piezo}
\end{equation}

where $x_\psi$ is the psi-induced displacement. For Si$_3$N$_4$ with
optical interferometric readout, displacement sensitivity is typically
$\sim 10^{-17}$~m$/\sqrt{\rm Hz}$. For AlN with direct piezoelectric
readout coupled to a SQUID or transmon:
\begin{equation}
x_{\rm min}^{\rm piezo} \sim \frac{\sqrt{4 k_B T / R_L}}{d_{33}/t}
\sim 10^{-15}~\text{m}/\sqrt{\text{Hz}} \quad (300~\text{K}),
\label{eq:x_min_piezo}
\end{equation}
which is $100\times$ worse than optical readout at room temperature.
At millikelvin temperatures with a transmon readout (quantum-limited):
\begin{equation}
x_{\rm min}^{\rm piezo,\,qm} \sim x_{\rm zpf}
= \sqrt{\frac{\hbar}{2 m_{\rm eff} \omega_m}}
\sim 10^{-16}~\text{m}.
\label{eq:x_zpf}
\end{equation}

\begin{finding}[AlN Piezoelectric PnC Assessment]
\label{find:aln}
\verdict{AlN PnC membranes do not compete with Si$_3$N$_4$ for
passive $Q$ performance ($120\times$ deficit). However, the
piezoelectric coupling to superconducting circuits provides a
hybrid readout pathway that eliminates optical interferometry
from the detection chain. Track as Phase~II readout alternative.
No change to primary material selection.}
\end{finding}


%=============================================================================
\section{Finding 5: Monocrystalline 4H-SiC Nanomechanical Resonators}
\label{sec:sic}
%=============================================================================

\subsection{Demonstrated Performance}

Sementilli et al.\ (Nano Letters, 2025) achieved:
\begin{itemize}[nosep]
\item Monocrystalline 4H-SiC nanostrings and trampolines
\item Damping as low as $\Gamma_{\rm min} = 2.7$~mHz at room temperature
\item $Q = 2\ee{7}$ at $\sim 100$~kHz (room temperature)
\item First observation of nonlinear dissipation in crystalline SiC
\item Material loss floor reached: $Q_{\rm int}^{-1} \sim 5\ee{-8}$
\item Low-stress operation (unlike Si$_3$N$_4$ which requires
  $\sigma > 1$~GPa)
\end{itemize}

\subsection{Path to $Q \sim 10^{10}$}

The authors project that combining their low-loss crystalline SiC
with dissipation dilution techniques should yield $Q \sim 10^{10}$
at room temperature. The argument:

\begin{enumerate}[nosep]
\item \textbf{Material $Q$ floor}: $Q_{\rm int} \sim 2\ee{7}$ (demonstrated).
\item \textbf{Dissipation dilution factor}: For Si$_3$N$_4$, soft
  clamping provides $D_Q \sim 10^3$--$10^4$ (from $Q \sim 10^5$
  bulk to $Q \sim 10^9$ soft-clamped).
\item \textbf{SiC projection}: $D_Q \sim 10^3$ applied to
  $Q_{\rm int} = 2\ee{7}$ gives $Q_{\rm eff} \sim 2\ee{10}$.
\end{enumerate}

This is \textit{consistent} with the crystalline Si result, where
$Q_{\rm int,Si} \sim 10^7$ at cryogenic $T$ + soft clamping gives
$Q > 10^{10}$.

\subsection{Unique SiC Advantages for DFD}

\begin{enumerate}[nosep]
\item \textbf{Room-temperature operation}: Unlike crystalline Si
  ($Q > 10^{10}$ only at 7~K), SiC could potentially achieve
  comparable $Q$ at 300~K, eliminating cryogenic requirements.

\item \textbf{Integrated spin-mechanical sensors}: Silicon vacancies
  (V$_{\rm Si}$) in 4H-SiC are analogous to NV centres in diamond
  but with better spectral stability and coherence at room
  temperature ($T_2 \sim 20$~ms). A SiC MEMS resonator with
  embedded V$_{\rm Si}$ defects could provide spin-based readout
  of mechanical displacement via strain coupling.

\item \textbf{Radiation hardness}: Relevant for any space-based
  deployment of DFD sensors (P6, P9).

\item \textbf{CMOS compatibility}: 4H-SiC wafers are commercially
  available (Wolfspeed, Coherent) and compatible with standard
  semiconductor processing.
\end{enumerate}

\begin{finding}[4H-SiC Nanomechanical Resonator Assessment]
\label{find:sic}
\honest{The projection of $Q \sim 10^{10}$ in SiC at room temperature
is physically plausible but \textit{entirely undemonstrated}. The
gap from demonstrated $Q = 2\ee{7}$ to projected $10^{10}$ is
$500\times$. While the Si$_3$N$_4$ precedent suggests this is
achievable ($10^4\times$ improvement via soft clamping demonstrated),
the SiC phononic crystal fabrication has not been attempted.}

\verdict{Add 4H-SiC to materials tracking list alongside Si$_3$N$_4$
and crystalline Si. Priority: LOW (behind both Si$_3$N$_4$ and Si).
Monitor for first demonstration of soft-clamped SiC PnC membrane.
If $Q > 10^9$ is achieved at room temperature, SiC becomes the
preferred Phase~III MEMS material (room-T + high-$Q$ + integrated
spin readout).}
\end{finding}


%=============================================================================
\section{Additional Material Assessments (Killed or Marginal)}
\label{sec:killed}
%=============================================================================

\subsection{Diamond NV Centres as psi-Field Sensors}

Diamond NV centres are exceptional magnetic field sensors
(sensitivity $\sim 1$~$\mu$T$/\sqrt{\rm Hz}$ on-chip, 2025) and
can couple to mechanical resonators via lattice strain ($g_{\rm sm}
\sim 2\pi \times 1$~MHz for diamond cantilevers). However, the DFD
psi-field couples to EM energy density $u_{\rm EM} = (E^2 + B^2)/2$,
not to magnetic field alone. The NV centre measures the Zeeman
splitting from $\mathbf{B}$, which is:

\begin{equation}
\Delta f_{\rm NV} = \gamma_{\rm NV} \cdot B
= 28~\text{GHz/T} \times B,
\label{eq:nv_zeeman}
\end{equation}

A psi-field perturbation $\dpsi$ modifies the local EM energy density
but does \textit{not} produce a net magnetic field --- the modification
is to the \textit{total} $u_{\rm EM}$, which is a scalar. An NV
centre measuring $|\mathbf{B}|$ would need to detect the second-order
effect of $\dpsi$ on the cavity field distribution, which is suppressed
by $\lam^2 \ll 1$.

\begin{finding}[Diamond NV Centre Assessment]
\label{find:nv}
\killed{Diamond NV centres as psi-field sensors: NV measures
$|\mathbf{B}|$ (vector), not $u_{\rm EM}$ (scalar). The psi-field
modifies $u_{\rm EM}$ without producing a net $\mathbf{B}$ perturbation
detectable by Zeeman splitting. Second-order effect suppressed by
$\lam^2 < 10^{-18}$ --- completely hopeless.}

\honest{NV centres remain useful as \textit{characterisation tools}
for SRF cavity magnetic field mapping (trapped flux diagnosis,
Abrikosov vortex imaging) but not as psi-sensors.}
\end{finding}

\subsection{Topological Insulator Surfaces}

Topological insulator (TI) surface states feature dissipationless
spin-momentum-locked Dirac fermions. The DFD question: could
TI surface conductivity provide an ultra-low-loss electromagnetic
resonator for passive detection?

The surface state conductivity is:
\begin{equation}
\sigma_{\rm TI} = \frac{e^2}{2h} \quad \text{(per surface, quantised)},
\label{eq:ti_conductivity}
\end{equation}
giving a surface resistance:
\begin{equation}
R_s^{\rm TI} = \frac{1}{\sigma_{\rm TI}}
= \frac{2h}{e^2} \approx 51.6~\text{k}\Omega.
\label{eq:ti_rs}
\end{equation}

This is \textit{enormous} compared to superconductor surface
resistance: $R_s^{\rm Nb} \sim 10$~n$\Omega$ at 20~mK. The TI
surface state has quantised \textit{DC} conductance, but at
microwave frequencies the relevant quantity is the AC surface
impedance, which includes both the quantum Hall contribution and
bulk conduction.

\begin{finding}[Topological Insulator Surface Assessment]
\label{find:ti}
\killed{Topological insulator surfaces for psi-detection:
$R_s^{\rm TI} \sim 50$~k$\Omega$ vs $R_s^{\rm Nb} \sim 10$~n$\Omega$
--- a factor of $5\ee{12}$ worse. TI surface conductivity is
quantised and dissipationless only for DC transport; at GHz
frequencies, losses from bulk carriers and phonon scattering dominate.
No advantage over superconducting surfaces for electromagnetic
cavity Q.}
\end{finding}

\subsection{Photonic Crystal Cavities for Optomechanical psi-Readout}

Photonic crystal (PhC) optomechanical cavities achieve optical $Q_{\rm opt}
\sim 10^6$ and mechanical $Q_m \sim 10^6$ simultaneously at GHz
frequencies, with optomechanical coupling $g_0/2\pi \sim 100$~kHz.
These are the basis of optomechanical crystals (OMCs) used in
quantum transduction experiments.

\textbf{DFD relevance}: OMCs could serve as the \textit{readout}
for MEMS psi-detection, coupling the mechanical displacement to
an optical cavity for quantum-limited measurement. However:
\begin{itemize}[nosep]
\item OMC mechanical $Q \sim 10^6$ is $1000\times$ below Si$_3$N$_4$
  PnC ($Q \sim 10^9$), so the OMC is not the \textit{sensor} ---
  it is the \textit{transducer}.
\item The readout noise floor is set by $\hbar\omega_{\rm opt}/(2
  Q_{\rm opt})$, giving displacement sensitivity
  $\sim 10^{-19}$~m$/\sqrt{\rm Hz}$ for $Q_{\rm opt} = 10^6$
  at 1550~nm.
\item This is better than standard Fabry-Perot interferometry
  ($\sim 10^{-17}$~m$/\sqrt{\rm Hz}$) but requires cryogenic
  operation (thermal occupation must be $\bar{n} < 1$).
\end{itemize}

\begin{finding}[Photonic Crystal OMC Assessment]
\label{find:phc}
\verdict{Photonic crystal OMCs are a viable readout technology
for MEMS psi-detection but do not change the sensor material
selection. Their $10^{-19}$~m$/\sqrt{\rm Hz}$ displacement
sensitivity is excellent for Phase~II/III, providing $100\times$
improvement over standard interferometry. Requires cryogenic
operation. Track as Phase~III readout upgrade.}
\end{finding}

\subsection{Metamaterials for psi-Field Enhancement}

Metamaterial-based electromagnetic energy density enhancement could
in principle amplify $u_{\rm EM}$ locally, increasing the psi-sourcing
strength. However, DFD couples to the \textit{volume-integrated}
$u_{\rm EM}$, and the Volume Cancellation Theorem (W3i5) proves
that active EM sourcing in a cavity produces psi-signals that
cancel over the cavity volume. Metamaterial local enhancement does
not circumvent this theorem --- it merely redistributes $u_{\rm EM}$
within the cavity without increasing the volume integral.

\begin{finding}[Metamaterial Enhancement Assessment]
\label{find:metamaterial}
\killed{Metamaterial EM energy density enhancement for psi-sourcing:
blocked by Volume Cancellation Theorem (W3i5). Local $u_{\rm EM}$
enhancement does not increase volume-integrated sourcing.
Furthermore, metamaterial losses ($Q \sim 10^2$--$10^4$) are
catastrophically worse than SRF cavities ($Q \sim 10^{10}$),
degrading passive FOM by $10^6$.}
\end{finding}


%=============================================================================
\section{Updated Materials Roadmap}
\label{sec:roadmap}
%=============================================================================

\begin{table}[H]
\centering
\caption{Updated materials roadmap for DFD detection programme (W4i13).}
\label{tab:roadmap}
\renewcommand{\arraystretch}{1.3}
\begin{tabular}{p{2cm}p{3cm}p{2.5cm}p{2cm}p{3cm}}
\toprule
\textbf{Phase} & \textbf{SRF Material} & \textbf{MEMS Material} &
\textbf{$T_{\rm op}$} & \textbf{Change from W4i10} \\
\midrule
I (2027--28) &
  N-doped Nb (primary); O-doped Nb (backup) &
  Si$_3$N$_4$ PnC (room $T$) &
  SRF: 20~mK; MEMS: 300~K &
  \textcolor{teal}{O-doped Nb added} \\
\midrule
II (2028--29) &
  N-doped Nb &
  Si$_3$N$_4$ PnC + crystalline Si PnC (cryo) &
  SRF: 20~mK; MEMS: 300~K / 7~K &
  \textcolor{teal}{Crystalline Si dual-track added} \\
\midrule
III (2030+) &
  Nb$_3$Sn &
  Best of: cm-scale Si$_3$N$_4$, cryo Si, SiC (if demonstrated) &
  SRF: 4~K; MEMS: TBD &
  \textcolor{teal}{Nb$_3$Sn at 4~K replaces Nb at 20~mK; SiC on
  watch list} \\
\bottomrule
\end{tabular}
\end{table}


%=============================================================================
\section{Consistency Checks and Flagged Issues}
\label{sec:consistency}
%=============================================================================

\subsection{Inconsistency Flag: Phase~I $Q_0$ Sensitivity Scaling}

W4i10 Eq.~(5) states $\lam_{\rm min} \propto 1/Q_0$, giving
degraded reach $\lam_{\rm min}^{\rm baseline} = 3.1\ee{-9}$ at
$Q_0 = 5\ee{10}$. However, the cooperativity $C \propto Q^2$
(from MASTER\_FINDINGS\_CORPUS \S2.4), which implies $\lam_{\rm min}
\propto 1/Q_0$ only in the $C \ll 1$ regime. At $Q_0 = 2\ee{11}$,
$C$ could approach unity, changing the scaling to $\lam_{\rm min}
\propto 1/Q_0^{1/2}$. This affects the Phase~I reach at target $Q_0$.

\caution{Verify whether Phase~I operates in $C \ll 1$ or $C \sim 1$
regime at $Q_0 = 2\ee{11}$. If $C \sim 1$, the reach improvement
from baseline to target is $\sqrt{4} = 2\times$, not $4\times$.}

\subsection{Inconsistency Flag: $Q_\psi$ Threshold Derivation}

The $\Qpsi$ threshold of $2.2\ee{8}$ (W4i11) was corrected from
$3\ee{9}$ (W4i10). The correction factor is $\sim 14\times$. This
should be cross-checked against the MEMS in-gap detection derivation
in W4i11\_Detection\_Update.tex to verify the phonon occupation
number at 300~K vs 7~K. At 7~K, the thermal occupation at 1~MHz is:
\begin{equation}
\bar{n}_{\rm th}(7~{\rm K}, 1~{\rm MHz})
= \frac{1}{e^{h f / k_B T} - 1}
\approx \frac{k_B T}{h f}
= \frac{1.38\ee{-23} \times 7}{6.63\ee{-34} \times 10^6}
\approx 1.46\ee{5},
\label{eq:nbar_cryo}
\end{equation}
vs at 300~K: $\bar{n}_{\rm th} \approx 6.25\ee{6}$. The ratio is
$43\times$, meaning the cryogenic Si resonator at 7~K has $43\times$
lower thermal noise, \textit{in addition to} the $10\times$ higher $Q$.
Combined improvement: $43 \times 10 = 430\times$ in SNR.

\confirmed{Cryogenic crystalline Si provides $\sim 430\times$
improvement in MEMS in-gap SNR over room-temperature Si$_3$N$_4$
(from $Q$ and thermal occupation combined).}


%=============================================================================
\section{Summary of Material Status Changes}
\label{sec:changes}
%=============================================================================

\begin{table}[H]
\centering
\caption{Material status changes from W4i10 to W4i13.}
\renewcommand{\arraystretch}{1.3}
\begin{tabular}{p{4cm}p{3cm}p{3cm}p{3cm}}
\toprule
\textbf{Material/Technology} & \textbf{W4i10 Status} &
\textbf{W4i13 Status} & \textbf{Reason} \\
\midrule
Crystalline Si PnC & Not assessed & \textbf{ADDED: Phase~II+} &
  $Q > 10^{10}$ at 7~K; $45\times$ margin \\
Nb$_3$Sn SRF & Not assessed & \textbf{ADDED: Phase~III} &
  $Q_0 \sim 10^{10}$ at 4~K; \$230M cryo savings \\
O-doped Nb & Not assessed & \textbf{ADDED: Phase~I backup} &
  SQMS March 2026; comparable $R_s$ to N-doping \\
AlN PnC & Not assessed & \textbf{WATCH: Readout option} &
  Piezo coupling; $Q$ too low for sensor \\
4H-SiC & Not assessed & \textbf{WATCH: Long-term} &
  $Q = 2\ee{7}$; path to $10^{10}$ at RT \\
cm-scale Si$_3$N$_4$ & Not assessed & \textbf{CONFIRMED: Phase~III} &
  $Q = 6.6\ee{9}$ at RT; $30\times$ margin \\
Levitated nanoparticles & Not assessed & \textbf{KILLED} &
  $250\times$ gap for grav.; $10^7\times$ for EM \\
Diamond NV centres & Not assessed & \textbf{KILLED} &
  Measures $\mathbf{B}$, not $u_{\rm EM}$; $\lam^2$ suppression \\
Topological insulators & Not assessed & \textbf{KILLED} &
  $R_s$ $5\ee{12}\times$ worse than Nb \\
Metamaterials & Not assessed & \textbf{KILLED} &
  Volume Cancellation Theorem; $Q$ too low \\
\midrule
Si$_3$N$_4$ PnC & Primary MEMS & \textbf{UNCHANGED: Primary} &
  Remains best near-term option \\
N-doped Nb & Primary SRF & \textbf{UNCHANGED: Primary} &
  Still best Phase~I material \\
He-3 & KILLED & \textbf{KILLED} & Confirmed dead \\
\bottomrule
\end{tabular}
\end{table}


%=============================================================================
\section{Recommendations for Next Iteration}
\label{sec:next}
%=============================================================================

\begin{enumerate}[nosep]
\item \textbf{Crystalline Si MEMS procurement}: Contact EPFL
  (Kippenberg/Ghadimi) or DTU Nanolab for collaboration on
  DFD-spec crystalline Si PnC resonators. Estimated 6--12 months
  to first device. Budget: \$100--200K within Phase~I.

\item \textbf{Nb$_3$Sn Phase~III design study}: Commission a
  conceptual design report for 256-cavity array at 4~K with
  Nb$_3$Sn. Engage Fermilab SRF group (already SQMS partner).
  Cost: \$50--100K for CDR.

\item \textbf{O-doped Nb recipe validation}: Request SQMS provide
  O-doped test cavities alongside N-doped for Phase~I.
  No additional cost if within SQMS programme.

\item \textbf{SiC soft-clamping feasibility}: Monitor literature
  for first demonstration of phononic crystal soft-clamping in
  monocrystalline SiC. If published, reassess priority.

\item \textbf{cm-scale Si$_3$N$_4$ trampoline}: Investigate whether
  the $Q = 6.6\ee{9}$ centimetre-scale trampoline resonator can be
  adapted for DFD in-gap detection. Key question: effective mass
  (likely $\gg 1.2$~ng, reducing FOM). Contact DTU (Schliesser group).
\end{enumerate}


\end{document}
